% ╔═══════════════════════════════════════╗
% ║   CHAPITRE 4 : LES SUITES            ║
% ╚═══════════════════════════════════════╝
\chapter{Les Suites Numériques}

\section{Définitions fondamentales}

\begin{defbox}[--- Suite numérique]
Une \textbf{suite} est une application $u : \N \to \R$ (ou $\C$). Pour $n \in \N$, on note $u(n) = u_n$ le \textbf{terme général}. La suite est notée $(u_n)_{n \in \N}$ ou $(u_n)_{n \geqslant n_0}$.
\end{defbox}

\begin{defbox}[--- Bornitude]
Soit $(u_n)$ une suite.
\begin{itemize}[leftmargin=1cm]
  \item $(u_n)$ est \textbf{majorée} si $\exists M \in \R,\; \forall n \in \N,\; u_n \leqslant M$.
  \item $(u_n)$ est \textbf{minorée} si $\exists m \in \R,\; \forall n \in \N,\; u_n \geqslant m$.
  \item $(u_n)$ est \textbf{bornée} si $\exists M \in \R,\; \forall n \in \N,\; |u_n| \leqslant M$.
\end{itemize}
\end{defbox}

\begin{defbox}[--- Monotonie]
\begin{itemize}[leftmargin=1cm]
  \item $(u_n)$ est \textbf{croissante} si $\forall n \in \N,\; u_{n+1} \geqslant u_n$, soit $u_{n+1} - u_n \geqslant 0$.
  \item Pour une suite à termes \textbf{strictement positifs} : $(u_n)$ croissante $\Leftrightarrow$ $\forall n,\; \frac{u_{n+1}}{u_n} \geqslant 1$.
\end{itemize}
\end{defbox}


\section{Limites}

\subsection{Définition formelle}

\begin{defbox}[--- Limite finie]
La suite $(u_n)$ a pour limite $\ell \in \R$ si :
\[
  \forall \varepsilon > 0,\; \exists N \in \N,\; \forall n \geqslant N,\; |u_n - \ell| \leqslant \varepsilon
\]
On note $\displaystyle\lim_{n \to +\infty} u_n = \ell$. La suite est alors dite \textbf{convergente}.
\end{defbox}

\begin{defbox}[--- Limite infinie]
$(u_n)$ tend vers $+\infty$ si :
\[
  \forall A > 0,\; \exists N \in \N,\; \forall n \geqslant N,\; u_n \geqslant A
\]
\end{defbox}

\begin{intuitionbox}[Lecture de la définition $\varepsilon$-$N$]
La définition dit : <<~quel que soit le niveau de précision $\varepsilon$ que l'on exige, il existe un rang $N$ à partir duquel tous les termes de la suite sont dans la bande $[\ell - \varepsilon,\, \ell + \varepsilon]$~>>.

\textbf{Attention à l'ordre :} $N$ \emph{dépend} de $\varepsilon$. Plus $\varepsilon$ est petit, plus $N$ est grand (en général). On ne peut \textbf{pas} échanger $\forall \varepsilon$ et $\exists N$.
\end{intuitionbox}

\subsection{Unicité de la limite}

\begin{propbox}[--- Unicité]
Si une suite converge, sa limite est unique.
\end{propbox}

\begin{proofbox}
Par l'absurde. Supposons $u_n \to \ell$ et $u_n \to \ell'$ avec $\ell \neq \ell'$. Posons $\varepsilon = \frac{|\ell - \ell'|}{2} > 0$.
\begin{itemize}[leftmargin=1cm]
  \item Il existe $N_1$ tel que $n \geqslant N_1 \Rightarrow |u_n - \ell| < \varepsilon$.
  \item Il existe $N_2$ tel que $n \geqslant N_2 \Rightarrow |u_n - \ell'| < \varepsilon$.
\end{itemize}
Pour $N = \max(N_1, N_2)$, l'inégalité triangulaire donne :
\[
  |\ell - \ell'| = |\ell - u_N + u_N - \ell'| \leqslant |u_N - \ell| + |u_N - \ell'| < 2\varepsilon = |\ell - \ell'|
\]
Contradiction : $|\ell - \ell'| < |\ell - \ell'|$. Donc $\ell = \ell'$. \hfill$\square$
\end{proofbox}

\subsection{Opérations sur les limites}

\begin{propbox}[--- Opérations (cas fini)]
Si $\displaystyle\lim u_n = \ell$ et $\displaystyle\lim v_n = \ell'$ (limites finies), alors :
\begin{enumerate}[leftmargin=1.5cm]
  \item $\displaystyle\lim (\lambda u_n) = \lambda \ell$ \quad pour $\lambda \in \R$.
  \item $\displaystyle\lim (u_n + v_n) = \ell + \ell'$.
  \item $\displaystyle\lim (u_n \cdot v_n) = \ell \cdot \ell'$.
  \item Si $\ell \neq 0$, alors $u_n \neq 0$ pour $n$ assez grand et $\displaystyle\lim \frac{1}{u_n} = \frac{1}{\ell}$.
\end{enumerate}
\end{propbox}

\begin{rembox}[Formes indéterminées]
Certaines combinaisons ne permettent pas de conclure directement :
\[
  +\infty - \infty, \qquad 0 \times \infty, \qquad \frac{\infty}{\infty}, \qquad \frac{0}{0}, \qquad 1^\infty
\]
Il faut alors une étude au cas par cas (factorisation, encadrement, etc.).
\end{rembox}


\section{Théorèmes d'encadrement}

\begin{thmbox}[--- Théorème des gendarmes]
Si $u_n \leqslant v_n \leqslant w_n$ pour tout $n$ et $\displaystyle\lim u_n = \lim w_n = \ell$, alors $(v_n)$ converge et $\displaystyle\lim v_n = \ell$.
\end{thmbox}

\begin{proofbox}
En soustrayant $(u_n)$, on se ramène à montrer : si $0 \leqslant v_n' \leqslant w_n'$ et $\lim w_n' = 0$, alors $\lim v_n' = 0$.

Soit $\varepsilon > 0$. Il existe $N$ tel que $n \geqslant N \Rightarrow |w_n'| < \varepsilon$. Comme $|v_n'| = v_n' \leqslant w_n' = |w_n'|$, on a $|v_n'| < \varepsilon$ pour $n \geqslant N$. \hfill$\square$
\end{proofbox}

\begin{propbox}[--- Suite convergente $\Rightarrow$ suite bornée]
Toute suite convergente est bornée.
\end{propbox}

\begin{proofbox}
Soit $(u_n)$ convergente vers $\ell$. Pour $\varepsilon = 1$, il existe $N$ tel que $n \geqslant N \Rightarrow |u_n| \leqslant |\ell| + 1$.

Posons $M = \max(|u_0|, |u_1|, \ldots, |u_{N-1}|, |\ell| + 1)$. Alors $|u_n| \leqslant M$ pour tout $n$. \hfill$\square$
\end{proofbox}

\begin{intuitionbox}[La réciproque est fausse !]
La suite $u_n = (-1)^n$ est bornée (par 1) mais diverge. La bornitude est nécessaire mais pas suffisante pour la convergence. Il faut ajouter la monotonie (théorème de la suite monotone) ou une autre condition.
\end{intuitionbox}


\section{Exemples remarquables}

\subsection{Suite géométrique}

\begin{propbox}[--- Suite géométrique $u_n = a^n$]
\begin{enumerate}[leftmargin=1.5cm]
  \item $a = 1$ : $u_n = 1$ pour tout $n$.
  \item $a > 1$ : $\lim u_n = +\infty$.
  \item $-1 < a < 1$ : $\lim u_n = 0$.
  \item $a \leqslant -1$ : $(u_n)$ diverge.
\end{enumerate}

\textbf{Somme partielle} ($a \neq 1$) :
\[
  \sum_{k=0}^{n} a^k = \frac{1 - a^{n+1}}{1 - a}
\]
Si $|a| < 1$ : $\displaystyle\sum_{k=0}^{+\infty} a^k = \frac{1}{1-a}$.
\end{propbox}

\begin{intuitionbox}[La somme télescopique cachée]
La preuve de la somme géométrique est un pattern télescopique : en multipliant $S = 1 + a + \cdots + a^n$ par $(1-a)$, les termes intermédiaires s'annulent deux à deux, ne laissant que $1 - a^{n+1}$.

Ce mécanisme revient constamment : somme télescopique, produit télescopique, ou même le calcul de l'inverse d'une matrice par élimination.
\end{intuitionbox}

\subsection{Le critère du rapport}

\begin{thmbox}[--- Critère $|u_{n+1}/u_n| < \ell < 1$]
Si $(u_n)$ est une suite de réels non nuls et s'il existe $\ell < 1$ tel que $\left|\frac{u_{n+1}}{u_n}\right| < \ell$ pour tout $n$ (ou à partir d'un certain rang), alors $\lim u_n = 0$.
\end{thmbox}

\begin{proofbox}
On écrit $\frac{u_n}{u_0} = \frac{u_1}{u_0} \cdot \frac{u_2}{u_1} \cdots \frac{u_n}{u_{n-1}}$, d'où $|u_n| < |u_0| \cdot \ell^n$. Comme $\ell < 1$, $\ell^n \to 0$ et donc $u_n \to 0$. \hfill$\square$
\end{proofbox}

\begin{exbox}[Application classique : $\lim \frac{a^n}{n!} = 0$]
Pour $u_n = \frac{a^n}{n!}$ (avec $a$ fixé), on a $\frac{u_{n+1}}{u_n} = \frac{a}{n+1} \to 0$. Donc $\lim u_n = 0$.

\textbf{Interprétation :} la factorielle croît plus vite que n'importe quelle exponentielle.
\end{exbox}


\section{Théorèmes de convergence}

\subsection{Suite monotone bornée}

\begin{thmbox}[--- Convergence monotone]
Toute suite croissante et majorée est convergente.

\medskip
\textit{Variantes :}
\begin{itemize}[leftmargin=1cm]
  \item Suite décroissante et minorée $\Rightarrow$ convergente.
  \item Suite croissante non majorée $\Rightarrow$ tend vers $+\infty$.
\end{itemize}
\end{thmbox}

\begin{proofbox}
Soit $A = \{u_n \mid n \in \N\}$. L'ensemble $A$ est non vide et majoré, donc admet une borne supérieure $\ell = \sup A$ (propriété de $\R$).

Soit $\varepsilon > 0$. Par caractérisation de la borne supérieure, il existe $u_N \in A$ tel que $\ell - \varepsilon < u_N \leqslant \ell$.

Comme $(u_n)$ est croissante, pour $n \geqslant N$ : $\ell - \varepsilon < u_N \leqslant u_n \leqslant \ell$, donc $|u_n - \ell| \leqslant \varepsilon$. \hfill$\square$
\end{proofbox}

\begin{intuitionbox}[Pourquoi ce théorème est si puissant]
Ce théorème permet de prouver la convergence \emph{sans connaître la limite}. On montre séparément :
\begin{enumerate}[leftmargin=1cm]
  \item La monotonie (souvent par récurrence ou par étude du signe de $u_{n+1} - u_n$).
  \item La bornitude (souvent par récurrence ou majoration).
\end{enumerate}
La limite est ensuite déterminée par passage à la limite dans la relation de récurrence (pour les suites récurrentes) ou par identification.
\end{intuitionbox}

\subsection{Suites adjacentes}

\begin{defbox}[--- Suites adjacentes]
Les suites $(u_n)$ et $(v_n)$ sont \textbf{adjacentes} si :
\begin{enumerate}[leftmargin=1cm]
  \item $(u_n)$ est croissante et $(v_n)$ est décroissante,
  \item $\forall n,\; u_n \leqslant v_n$,
  \item $\lim(v_n - u_n) = 0$.
\end{enumerate}
\end{defbox}

\begin{thmbox}[--- Convergence des suites adjacentes]
Si $(u_n)$ et $(v_n)$ sont adjacentes, elles convergent vers la même limite $\ell$, avec :
\[
  u_0 \leqslant u_1 \leqslant \cdots \leqslant u_n \leqslant \cdots \leqslant \ell \leqslant \cdots \leqslant v_n \leqslant \cdots \leqslant v_1 \leqslant v_0
\]
\end{thmbox}

\begin{proofbox}
$(u_n)$ est croissante et majorée par $v_0$ : elle converge vers $\ell$.

$(v_n)$ est décroissante et minorée par $u_0$ : elle converge vers $\ell'$.

Enfin $\ell' - \ell = \lim(v_n - u_n) = 0$, donc $\ell' = \ell$. \hfill$\square$
\end{proofbox}


\subsection{Suites extraites et Bolzano-Weierstrass}

\begin{defbox}[--- Suite extraite]
Une \textbf{suite extraite} de $(u_n)$ est une suite $(u_{\varphi(n)})$ où $\varphi : \N \to \N$ est strictement croissante.
\end{defbox}

\begin{propbox}
Si $\lim u_n = \ell$, alors toute suite extraite converge aussi vers $\ell$.

\textbf{Contraposée :} si deux sous-suites convergent vers des limites distinctes, la suite diverge.
\end{propbox}

\begin{thmbox}[--- Bolzano-Weierstrass]
Toute suite bornée admet une sous-suite convergente.
\end{thmbox}

\begin{proofbox}[Démonstration (par dichotomie)]
Les valeurs de $(u_n)$ sont contenues dans $[a_0, b_0]$. On coupe en deux : au moins l'un des sous-intervalles contient $u_n$ pour une infinité d'indices. On note $[a_1, b_1]$ cet intervalle et on choisit $\varphi(1) > \varphi(0)$ tel que $u_{\varphi(1)} \in [a_1, b_1]$.

En itérant, on construit :
\begin{itemize}[leftmargin=1cm]
  \item des intervalles emboîtés $[a_k, b_k]$ de longueur $\frac{b_0 - a_0}{2^k} \to 0$,
  \item une suite extraite $(u_{\varphi(n)})$ avec $u_{\varphi(n)} \in [a_n, b_n]$.
\end{itemize}
Les suites $(a_n)$ et $(b_n)$ sont adjacentes, elles convergent vers un même $\ell$. Par le théorème des gendarmes, $u_{\varphi(n)} \to \ell$. \hfill$\square$
\end{proofbox}


\section{Suites récurrentes}

\begin{defbox}[--- Suite récurrente]
Soit $f : \R \to \R$. Une suite récurrente est définie par :
\[
  u_0 \in \R \quad \text{et} \quad u_{n+1} = f(u_n) \;\text{ pour } n \geqslant 0.
\]
\end{defbox}

\begin{propbox}[--- Points fixes et limite]
Si $f$ est continue et $(u_n)$ converge vers $\ell$, alors $f(\ell) = \ell$ \quad ($\ell$ est un \textbf{point fixe} de $f$).
\end{propbox}

\begin{proofbox}
Comme $u_n \to \ell$ et $f$ continue, $f(u_n) \to f(\ell)$. Or $f(u_n) = u_{n+1} \to \ell$. Par unicité de la limite, $f(\ell) = \ell$. \hfill$\square$
\end{proofbox}

\begin{methodbox}[Étude complète d'une suite récurrente $u_{n+1} = f(u_n)$]
\begin{enumerate}[leftmargin=1.5cm]
  \item \textbf{Tracer} le graphe de $f$ et la droite $y = x$ (indispensable).
  \item \textbf{Points fixes :} résoudre $f(\ell) = \ell$.
  \item \textbf{Stabilité :} montrer que $f([a,b]) \subset [a,b]$ pour un intervalle bien choisi.
  \item \textbf{Monotonie :} étudier le signe de $f(x) - x$ ou utiliser la croissance/décroissance de $f$.
  \item \textbf{Convergence :} appliquer le théorème de convergence monotone.
  \item \textbf{Identifier la limite :} parmi les points fixes.
\end{enumerate}
\end{methodbox}

\subsection{Cas $f$ croissante}

\begin{propbox}
Si $f : [a,b] \to [a,b]$ est continue et croissante, et $u_0 \in [a,b]$, alors $(u_n)$ est \textbf{monotone} et converge vers un point fixe $\ell \in [a,b]$.
\begin{itemize}[leftmargin=1cm]
  \item Si $u_1 \geqslant u_0$ : $(u_n)$ croissante.
  \item Si $u_1 \leqslant u_0$ : $(u_n)$ décroissante.
\end{itemize}
\end{propbox}

\subsection{Cas $f$ décroissante}

\begin{propbox}
Si $f : [a,b] \to [a,b]$ est continue et décroissante, alors :
\begin{itemize}[leftmargin=1cm]
  \item La sous-suite $(u_{2n})$ converge vers $\ell$ vérifiant $f \circ f(\ell) = \ell$.
  \item La sous-suite $(u_{2n+1})$ converge vers $\ell'$ vérifiant $f \circ f(\ell') = \ell'$.
\end{itemize}
On étudie $f \circ f$ (qui est croissante) pour se ramener au cas précédent.
\end{propbox}

\begin{exbox}[Le nombre d'or comme point fixe]
Soit $f(x) = 1 + \frac{1}{x}$ et $u_0 > 0$ avec $u_{n+1} = f(u_n)$.

$f$ est décroissante sur $]0, +\infty[$. On calcule $f \circ f(x) = \frac{2x+1}{x+1}$.

Les points fixes de $f \circ f$ : $\frac{2x+1}{x+1} = x \Leftrightarrow x^2 - x - 1 = 0$, dont la racine positive est $\ell = \frac{1 + \sqrt{5}}{2}$ --- le nombre d'or.

On montre que $(u_{2n})$ et $(u_{2n+1})$ convergent toutes deux vers $\varphi = \frac{1+\sqrt{5}}{2}$.
\end{exbox}

\begin{intuitionbox}[Escalier vs escargot]
\textbf{$f$ croissante :} le graphe donne un <<~escalier~>> (monotone). Les termes montent ou descendent vers le point fixe.

\textbf{$f$ décroissante :} le graphe donne un <<~escargot~>> (spirale). Les termes alternent de part et d'autre du point fixe, se rapprochant en spirale.

Toujours tracer le graphe : il dit immédiatement si la suite converge et vers quel point fixe.
\end{intuitionbox}


\section{Synthèse personnelle --- Suites}

\begin{intuitionbox}[Connexions identifiées]
\textbf{1. La logique au service de l'analyse.} La définition $\varepsilon$-$N$ de la convergence est une phrase logique avec quantificateurs imbriqués. Nier cette définition (pour montrer la divergence) exige exactement les compétences du chapitre 1.

\textbf{2. Les complexes dans les suites.} La somme géométrique $\sum a^k = \frac{1 - a^{n+1}}{1-a}$ est valable pour $a \in \C$. Les racines de l'unité (chapitre 2) sont reliées aux suites récurrentes via les suites définies par $u_{n+1} = e^{2i\pi/n} \cdot u_n$.

\textbf{3. Pattern récurrent : séparer existence et calcul.} Les théorèmes de convergence monotone et de Bolzano-Weierstrass prouvent l'\emph{existence} d'une limite sans la calculer. L'identification de la limite (point fixe, borne sup) est une étape distincte. Cette séparation est un principe fondamental en mathématiques.

\textbf{4. La borne supérieure comme clé de voûte.} Le théorème de convergence monotone repose sur la propriété de la borne supérieure de $\R$, qui est l'axiome fondamental distinguant $\R$ de $\Q$. C'est pourquoi les suites de rationnels peuvent converger vers un irrationnel (par ex. les approximations décimales de $\pi$).
\end{intuitionbox}


